\documentclass[12pt]{article}
\usepackage[utf8]{inputenc}
\usepackage[T1]{fontenc}
\usepackage{hyperref}
\title{Lower Pokomo Language Resources for K-12 Curriculum Development}
\author{Compiled by Claude Sonnet 5.5 (claude-sonnet-5-5) for Andr\'as Kornai}
\date{}
\begin{document}
\maketitle

\section{Introduction}
Lower Pokomo is spoken along the lower Tana River in Kenya (Tana Delta, Garsen, Kipini area). The ISO 639-3 code \texttt{poj} has been retired: it was merged into Kipfokomo (Pokomo) \texttt{pkb} (change request 2007-170, effective 2008). Glottolog also notes that the Upper/Lower Pokomo split is not primarily one of intelligibility. All resources found in search are therefore filed under Kipfokomo/Pokomo; none was found that is specific to the Lower variety alone, except where noted. Sources disagree on speaker numbers (roughly 29,000--34,000 in some listings, about 95,000 in a 2009 census figure cited elsewhere). Little is available for school use.

\section{Primary Educational Resources}
No primers, textbooks or school curricula for Pokomo were found in the search results.

\section{Linguistic and Reference Materials}
\begin{itemize}
\item Glottolog entry for the retired code poj (Lower Pokomo): \url{https://glottolog.org/resource/languoid/iso/poj}
\item Wikipedia, Pokomo language (background and code history): \url{https://en.wikipedia.org/wiki/Pokomo_language}
\end{itemize}
No dictionary or grammar of Pokomo appeared in the search results.

\section{Religious and Community Resources}
\begin{itemize}
\item Scripture Earth page for Pokomo (pkb): a New Testament (Ḍamano Ipfya) on Bible.com, a Bible app and a Gospel of Mark film are linked: \url{https://www.scriptureearth.org/00eng.php?iso=pkb}
\item Joshua Project, Lower Pokomo in Kenya (people-group information): \url{https://joshuaproject.net/people_groups/13120}
\end{itemize}

\section{Digital and Audio Resources}
\begin{itemize}
\item Global Recordings Network, Kipfokomo: audio and Bible-story recordings: \url{https://globalrecordings.net/en/language/pkb}
\end{itemize}

\section{Curriculum Development Context}
Evidence of Pokomo in schools was not found. Scripture and evangelistic audio exist, but a K-12 curriculum would need new materials (literacy primers, a dictionary, a grammar sketch), ideally developed with Pokomo communities and Kenyan education authorities. This dossier is deliberately short because little more was found.
\end{document}